% Source: https://artofproblemsolving.com/community/c3h2592654_ssa_congruence
% Stable identity: aops-topic-2592654-q01
Let $a=BC$. By the law of cosines, a triangle with the given data exists with $BC=a$ exactly when $a>0$ and
\[a^2-(2b\cos\gamma)\,a+\left(b^2-c^2\right)=0.\tag{$*$}\]
Triangles with the same $a$, $b$, $\gamma$ are congruent (SAS), so the number of triangles is the number of positive roots of $(*)$. Its roots have sum $2b\cos\gamma$, product $b^2-c^2$ and discriminant $4\left(c^2-b^2\sin^2\gamma\right)$.

\textbf{(b)}
\begin{itemize}
\item $c>b$: the product is negative, so there is exactly one positive root. One triangle.
\item $c=b$: the roots are $0$ and $2b\cos\gamma$, so there is one triangle if $\gamma<90^\circ$ and none otherwise.
\item $c<b$: the product is positive, so the roots have the same sign (or are not real). If $\gamma\ge90^\circ$ the sum is $\le0$, so there is no positive root. If $\gamma<90^\circ$ the sum is positive, and:
\begin{itemize}
\item for $c<b\sin\gamma$ there are no real roots (no triangle);
\item for $c=b\sin\gamma$ there is a double root (one triangle, right-angled at $B$);
\item for $b\sin\gamma<c<b$ there are two positive roots (two triangles, the ``ambiguous case'').
\end{itemize}
\end{itemize}
Collecting the cases with exactly one triangle gives the stated condition. $\blacksquare$

\textbf{(a)} Here $c=\sqrt{229}>12=b$, so by (b) the triangle is unique. Indeed, $(*)$ reads $a^2-12a+144-229=0$, i.e. $(a-17)(a+5)=0$, and the only positive root is $BC=\boxed{17}$.

\begin{asy}
usepackage("amsmath");
size(5.5cm);
defaultpen(fontsize(10pt));
pair C=(0,0);real g=50;real bb=4;real cc=3.3;
pair A=bb*dir(g);
real h=A.y; real dx=sqrt(cc^2-h^2);
pair B1=(A.x-dx,0),B2=(A.x+dx,0);
draw(C--(6.4,0),linewidth(0.7));draw(C--A,linewidth(0.8));
draw(A--B1,linewidth(0.8));draw(A--B2,linewidth(0.8));
draw(arc(A,cc,-118,-62),dashed);
draw((A.x,0)--A,dotted);
dot(A);dot(B1);dot(B2);
label("$C$",C,SW);label("$A$",A,N);label("$B_1$",B1,S);label("$B_2$",B2,S);
label("$b$",A/2,NW);label("$c$",(A+B1)/2,W,fontsize(9pt));label("$c$",(A+B2)/2,E,fontsize(9pt));
label("$\gamma$",C+(0.6,0.25),fontsize(9pt));
\end{asy}

The ambiguous case $b\sin\gamma<c<b$: the circle of radius $c$ about $A$ meets the other ray from $C$ at two points $B_1$ and $B_2$.

\begin{asy}
usepackage("amsmath");
size(5.5cm);
defaultpen(fontsize(10pt));
pair C=(0,0);pair A=12*dir(60);real a=17;pair B=(a,0);
draw(C--B--A--cycle,linewidth(0.8));
draw((-6,0)--(18.5,0),gray);
pair Bn=(-5,0);
draw(arc(A,sqrt(229),-148,-127),dashed);draw(arc(A,sqrt(229),-55,-33),dashed);
dot(Bn);label("$B'$",Bn,SW);
label("$C$",C,S);label("$A$",A,N);label("$B$",B,S);
label("$12$",A/2,W);label("$\sqrt{229}$",(A+B)/2,NE);label("$17$",(B+C)/2,S);
label("$60^\circ$",C+(1.6,0.6),fontsize(9pt));
\end{asy}

Part (a): the circle of radius $\sqrt{229}$ about $A$ meets line $CB$ at $B$, with $CB=17$, and at a point $B'$ on the wrong side of $C$, with $CB'=5$; this point corresponds to the negative root $-5$.
